\section{Revisión de contenidos}

\subsection{Preguntas sobre Concurrencia y Exclusión Mutua}

\begin{enumerate}
  \item Enumere cuatro problemas de diseño para los cuales el concepto de concurrencia es relevante.
  \item ¿Cuáles son tres contextos en los que surge la concurrencia?
  \item ¿Qué es una condición de carrera?
  \item Enumere tres grados de conciencia entre procesos y defina brevemente cada uno.
  \item ¿Cuál es la distinción entre procesos competidores (que compiten) y procesos cooperadores (que cooperan)?
  \item ¿Qué es la inanición con respecto al control de concurrencia mediante exclusión mutua?
  \item Enumere y describa brevemente el funcionamiento de las tres instrucciones de hardware para solucionar el problema de la exclusión mutua.
  \item ¿Qué es la espera activa o \emph{busy waiting}?
  \item ¿Qué operaciones se pueden realizar sobre un semáforo?
  \item ¿Cuál es la diferencia entre semáforos binarios y semáforos generales?
  \item ¿Cuál es la diferencia entre un mutex y un semáforo binario?
  \item ¿Qué características de los monitores los marcan como herramientas de sincronización de alto nivel?
  \item Compare el direccionamiento directo e indirecto con respecto al paso de mensajes.
\end{enumerate}

\subsection{Preguntas sobre interbloqueo}

\begin{enumerate}
  \item De ejemplos de recursos reutilizables y recursos consumibles.
  \item ¿Cuáles son las tres condiciones que deben estar presentes para que un interbloqueo sea posible?
  \item ¿Cuáles son las cuatro condiciones que provocan un interbloqueo?
  \item ¿Cómo puede prevenirse la condición de retención y espera?
  \item ¿Por qué no puede desactivarse la exclusión mutua para prevenir interbloqueos?
  \item ¿Cómo se puede predecir la espera circular?
  \item Explique brevemente el algoritmo del banquero.
  \item Liste algunos métodos que permiten recuperación tras un interbloqueo.
\end{enumerate}
